\section*{Hoofdstuk \ref{ch:g10:algebra} --- Algebra: vergelijkingen en ongelijkheden}

\begin{solution}{exo:g10:algebra:1}
$3(2x-5) - 2(x-7) = 6x - 15 - 2x + 14 = 4x - 1$.

$(x+4)(2x-1) = 2x^2 - x + 8x - 4 = 2x^2 + 7x - 4$.

$(3x-2)^2 = 9x^2 - 12x + 4$.

$(5-x)(5+x) = 25 - x^2$.
\end{solution}

\begin{solution}{exo:g10:algebra:2}
$x^2 - 49 = (x+7)(x-7)$ (verschil van kwadraten).

$x^2 + 10x + 25 = (x+5)^2$.

$7x^2 - 14x = 7x(x - 2)$ (gemeenschappelijke factor).

$16x^2 - 8x + 1 = (4x - 1)^2$.
\end{solution}

\begin{solution}{exo:g10:algebra:3}
$4x + 9 = x - 3$: $3x = -12$, dus $x = -4$.

$\frac{2x-1}{3} = 5$: $2x - 1 = 15$, dus $x = 8$.

$2(x-3) = 2x + 1$ werkt uit tot $2x - 6 = 2x + 1$, dus $-6 = 1$: onwaar voor
elke $x$. Geen oplossing --- de twee leden beschrijven evenwijdige rechten
die elkaar nooit ontmoeten.
\end{solution}

\begin{solution}{exo:g10:algebra:4}
$(x+5)(2x-8) = 0$: $x = -5$ of $x = 4$.

$x^2 - 16 = (x-4)(x+4) = 0$: $x = 4$ of $x = -4$.

$x^2 + 6x + 9 = (x+3)^2 = 0$: de enkele oplossing $x = -3$.
\end{solution}

\begin{solution}{exo:g10:algebra:5}
$3x - 5 \leq 7$: $3x \leq 12$, dus $x \leq 4$: oplossingsverzameling
$\intoc{-\infty}{4}$.

$-4x + 1 > 9$: $-4x > 8$, en delen door $-4$ keert de ongelijkheid om:
$x < -2$: oplossingsverzameling $\intoo{-\infty}{-2}$.

$2(x-1) \geq 5x + 4$: $2x - 2 \geq 5x + 4$, dus $-3x \geq 6$, dus
$x \leq -2$: oplossingsverzameling $\intoc{-\infty}{-2}$.
\end{solution}

\begin{solution}{exo:g10:algebra:6}
De gemeenschappelijke factor is $(x-1)$:
\[
E = (x-1)\bigl[(x+3) - (2x - 5)\bigr] = (x-1)(x + 3 - 2x + 5)
= (x-1)(8 - x).
\]
Dan is $E = 0$ voor $x = 1$ of $x = 8$ (nulproductregel).
\end{solution}

\begin{solution}{exo:g10:algebra:7}
$x^2 = 5x$: $x^2 - 5x = x(x - 5) = 0$, dus $x = 0$ of $x = 5$ (deel nooit
door $x$: de oplossing $0$ zou verloren gaan).

$(2x+1)^2 = (x-4)^2$: breng naar links en ontbind het verschil van
kwadraten:
\[
(2x+1)^2 - (x-4)^2
= \bigl[(2x+1) + (x-4)\bigr]\bigl[(2x+1) - (x-4)\bigr]
= (3x - 3)(x + 5).
\]
Nulproductregel: $x = 1$ of $x = -5$.
\end{solution}

\begin{solution}{exo:g10:algebra:8}
$(x+2)(4-x) > 0$: de factor $x + 2$ verdwijnt in $-2$ (daarna positief),
$4 - x$ verdwijnt in $4$ (daarvóór positief). Het product is positief precies
wanneer beide factoren positief zijn, dus op $\intoo{-2}{4}$.

$\frac{2x-6}{x+1} \geq 0$: teller nul in $3$, noemer nul in $-1$ (verboden).
\omterm{def:g10:algebra:signtable}{Tekentabel}: het quotiënt is positief wanneer beide delen hetzelfde teken
hebben, dus voor $x < -1$ of $x > 3$; het verdwijnt in $x = 3$.
Oplossingsverzameling: $\intoo{-\infty}{-1} \cup \intco{3}{+\infty}$.
\end{solution}

\begin{solution}{exo:g10:algebra:9}
Na $x$ maanden kost de eerste dienst $9x + 15$ en de tweede $12x$. De eerste
is goedkoper wanneer
\[
9x + 15 < 12x
\ \Longleftrightarrow\
15 < 3x
\ \Longleftrightarrow\
x > 5 .
\]
Vanaf de $6$de maand is de eerste dienst in totaal de goedkoopste keuze.
\end{solution}

\begin{solution}{exo:g10:algebra:10}
Verboden waarde: $x = 1$. Voor $x \neq 1$ vermenigvuldig je beide leden met
$x - 1$:
\[
x + 3 = 2(x - 1) = 2x - 2
\ \Longleftrightarrow\
x = 5 .
\]
Omdat $5 \neq 1$ is de oplossing $x = 5$. Controle:
$\frac{5+3}{5-1} = \frac84 = 2$.
\end{solution}

\begin{solution}{exo:g10:algebra:11}
$(n+1)^2 - n^2 = n^2 + 2n + 1 - n^2 = 2n + 1$, en dat is oneven voor elk
\omterm{def:g10:numbers:sets}{geheel getal} $n$.

Twee opeenvolgende oneven getallen schrijf je als $2n - 1$ en $2n + 1$. Dan
is
\[
(2n+1)^2 - (2n-1)^2
= \bigl[(2n+1)+(2n-1)\bigr]\bigl[(2n+1)-(2n-1)\bigr]
= 4n \times 2 = 8n,
\]
een veelvoud van $8$.
\end{solution}

\begin{solution}{pb:g10:algebra:1}
\textbf{1.} $x^2 - 10x + 25 = (x - 5)^2$;
$9x^2 - 49 = (3x - 7)(3x + 7)$;
$x^2 + 6x + 9 - 4 = (x + 3)^2 - 2^2 = (x + 1)(x + 5)$.

\textbf{2.} $x = -1$ of $x = -5$; \quad $x = \frac73$ of $x = -\frac73$.

\textbf{3.} $(5 - t)(5 + t) = 25 - t^2 = 21$ geeft $t^2 = 4$, dus $t = 2$: de
getallen zijn $3$ en $7$ (som $10$, product $21$: klopt).

\textbf{4.} $(7 - t)(7 + t) = 49 - t^2 = 33$: $t = 4$, getallen $3$ en $11$.

\textbf{5.} $\left(\frac s2 - t\right)\left(\frac s2 + t\right)
= \left(\frac s2\right)^2 - t^2 = p$ dwingt
$t^2 = \left(\frac s2\right)^2 - p$ af. Er is een oplossing precies wanneer
dit $\geq 0$ is, dus wanneer $p \leq \left(\frac s2\right)^2$ --- het product
van twee getallen met gegeven som kan het kwadraat van de halve som nooit
overtreffen (de weekendopgave over de omheining van koningin Dido in het
onderbouwvolume wist dat al).

\textbf{6.} $(x - a)(x - b) = x^2 - (a + b)x + ab$: de \omterm{def:g10:algebra:equation}{vergelijking}
$x^2 - sx + p = 0$ zegt precies ``$x$ is een van twee getallen met som $s$ en
product $p$''.

\textbf{7.} $(x-5)^2 - 4 = x^2 - 10x + 25 - 4 = x^2 - 10x + 21$. Verschil van
kwadraten: $(x - 5 - 2)(x - 5 + 2) = (x - 7)(x - 3) = 0$: oplossingen $3$ en
$7$ --- de getallen van de schrijver uit vraag 3.

\textbf{8.} $x^2 + 6x - 7 = (x + 3)^2 - 16 = 0$ geeft $(x + 3)^2 = 16$, dus
$x + 3 = \pm 4$: $x = 1$ of $x = -7$.

\textbf{9.} $x^2 + 4x + 5 = (x + 2)^2 + 1 \geq 1 > 0$ voor elke $x$: geen
reële oplossing. In de vorm $(x + h)^2 + k$ bestaan er oplossingen precies
wanneer $k \leq 0$ (dan heeft $(x+h)^2 = -k$ de twee wortels
$-h \pm \sqrt{-k}$); is $k > 0$, dan blijft de uitdrukking positief.

\textbf{10.} Het vierkant met zijde $x$ (oppervlakte $x^2$) met twee
rechthoeken $3 \times x$ ($6x$ samen) vormt een L met oppervlakte
$x^2 + 6x = 7$; de ontbrekende hoek $3 \times 3$ (oppervlakte $9$) vult het
aan tot een groot vierkant met zijde $x + 3$. $9$ bij beide leden optellen
geeft $(x + 3)^2 = 16$, dus $x + 3 = 4$ (lengten zijn positief) en $x = 1$.
De algebra en de tekening zijn dezelfde handeling.

\textbf{11.} $(x-3)(x+2)$: negatief tussen de nulpunten. Oplossing van
$< 0$: $x \in \intoo{-2}{3}$.

\textbf{12.} $x^2 - 10x + 21 \leq 0$, dus $(x - 3)(x - 7) \leq 0$: tussen de
nulpunten, grenzen meegerekend: $x \in \intcc{3}{7}$.

\textbf{13.} Nul in $x = 1$, verboden in $x = -2$; het quotiënt is positief
buiten de nulpunten: $x \in \intoo{-\infty}{-2} \cup \intco{1}{+\infty}$.

\textbf{14.} $25 - (x - 5)^2 = 25 - x^2 + 10x - 25 = x(10 - x)$: de
identiteit klopt. Omdat $(x-5)^2 \geq 0$ altijd geldt, is $A(x) \leq 25$, met
gelijkheid precies in $x = 5$: de vierkante wei $5 \times 5$, met maximum
$25$ m$^2$ --- de tabel uit jaar 6, nu een bewijs van twee regels.

\textbf{15.} Vermenigvuldig de bewering voor $x > 0$ met $x$:
$x^2 + 1 \geq 2x$, dus $x^2 - 2x + 1 = (x - 1)^2 \geq 0$ --- waar, met
gelijkheid alleen in $x = 1$. Terugdelen door $x > 0$ behoudt alles
(\cref{prop:g10:algebra:ineqrules}).

\textbf{16.} Kruiselings vermenigvuldigen in $\frac1x = \frac{x}{1 - x}$
(alle lengten positief): $x^2 = 1 - x$, dus $x^2 + x - 1 = 0$.

\textbf{17.} $x^2 + x - 1 = \left(x + \frac12\right)^2 - \frac54 = 0$ geeft
$x + \frac12 = \pm\frac{\sqrt5}{2}$. De positieve wortel:
$x = \frac{\sqrt5 - 1}{2} \approx 0.618$.

\textbf{18.} $\varphi = \frac{2}{\sqrt5 - 1} =
\frac{2(\sqrt5 + 1)}{(\sqrt5 - 1)(\sqrt5 + 1)}
= \frac{2(\sqrt5 + 1)}{5 - 1} = \frac{\sqrt5 + 1}{2} \approx 1.618$.

\textbf{19.} $\varphi^2 = \frac{(\sqrt5 + 1)^2}{4} = \frac{6 + 2\sqrt5}{4}
= \frac{3 + \sqrt5}{2}$, en $\varphi + 1 = \frac{\sqrt5 + 3}{2}$: gelijk. En
$\varphi - 1 = \frac{\sqrt5 - 1}{2} = x = \frac1\varphi$ (zo werd $\varphi$
juist gedefinieerd). Het gulden getal is het getal waarvan het kwadraat er
één bij optelt en het omgekeerde er één van aftrekt.

\textbf{20.} $\sqrt5$ is \omterm{ex:g10:numbers:classify}{irrationaal} (de weekendopgave over irrationaliteit
in het onderbouwvolume: $5$ is geen volkomen kwadraat), en $1$ optellen en
daarna halveren behoudt de irrationaliteit (rationale bewerkingen,
\cref{pb:g10:numbers:1}). Verhoudingen: $\frac53 \approx 1.667$,
$\frac85 = 1.600$, $\frac{13}{8} = 1.625$, $\frac{21}{13} \approx 1.615$:
afwisselend boven en onder, en steeds dichter bij
$\varphi = 1.618\ldots$ --- de ritmegetallen uit de weekendopgave over het
tellen van ritmes in het onderbouwvolume jagen in het geheim de gulden snede
na; de hoofdstukken over rijen zullen ze op heterdaad betrappen.
\end{solution}
